\documentclass[11pt]{article}
\usepackage[a4paper,margin=1in]{geometry}
\usepackage{amsmath,amssymb,amsthm,mathtools}
\usepackage{enumitem}
\usepackage[colorlinks=true,linkcolor=blue,citecolor=blue,urlcolor=blue]{hyperref}
\usepackage{xurl}

\newcommand{\indep}{\mathrel{\perp\mkern-9mu\perp}}
\newcommand{\Renc}{\mathcal{R}_{\mathrm{enc}}}
\newcommand{\epscpa}{\varepsilon_{\mathrm{cpa}}}
\newcommand{\HunpC}{H_{\mathrm{unp}}^{\mathcal{C}}}
\newcommand{\Hunp}{H_{\mathrm{unp}}}
\newcommand{\C}{\mathcal{C}}
\newcommand{\Yspace}{\mathcal{Y}}
\newcommand{\Wspace}{\mathcal{W}}
\newcommand{\Encf}{\mathrm{Enc}}
\newcommand{\Decf}{\mathrm{Dec}}
\newcommand{\AliceView}{\mathit{AliceView}}
\newcommand{\AliceViewHyb}{\mathit{AliceView}_{\textrm{after-swaps}}'}
\newcommand{\Dk}{\mathit{Dk}}
\newcommand{\pk}{\mathit{pk}}

\title{3-party DSDP $V_2$-unpredictability in $\AliceView$ against corrupted Alice}
\author{}
\date{}

\begin{document}
\maketitle

\section*{Setup}

\begin{enumerate}
\item $(\Omega, \Sigma, P)$ is a discrete probability space.

\item $(R, +, \cdot, 0_R, 1_R)$ is a finite commutative ring with $m := \#|R| \geq 2$. The unit group is $R^\times := \{ u \in R \mid \exists v \in R,\ u \cdot v = 1_R \}$.

\item $\Encf$ is the encryption procedure of a public-key encryption scheme with plaintext space $R$ and finite encryption-randomness space $\Renc$, and $\Decf$ is the corresponding decryption procedure. The fixed public-key registry $(\pk_a, \pk_b, \pk_c)$ and corresponding fixed decryption keys $(\Dk_a, \Dk_b, \Dk_c)$ are deterministic constants assumed to satisfy the real-or-zero bound stated in Setup item~8. The key-generation experiment is not formalized in this note. Fresh encryption randomness is sampled from $\Renc$ uniformly, i.e.\ each element of $\Renc$ is drawn with probability $1 / |\Renc|$. Decryption correctness $\Decf(\Dk_p, \Encf(\pk_p, m; r)) = m$ holds for every party $p$, message $m$, and randomness $r$; this identity is invoked implicitly in Step~3 (own-key decryption).

\item Each decryption key $\Dk_p$ is a deterministic value fixed at protocol setup and held privately by party $p$. When party $p$ is the corrupted party in a security theorem, $\Dk_p$ appears inside $p$'s own view since $p$ owns it. In such a case, the ciphertext encrypted by $\pk_p$ in party $p$'s view is a plaintext by decryption.

\item Public keys $(\pk_a, \pk_b, \pk_c)$ are known to all parties.

\item \textbf{The adversary class $\C$.} The corrupted party in this theorem is Alice, and the party-level adversary is static semi-honest: Alice follows the protocol, but an allowed predictor may use the whole corrupted-Alice view to guess $V_2$. This predictor is denoted $A$ below. It receives only a party view value, such as $\AliceView$, $H_1$, or $H_2$; it does not receive the reduction-side auxiliary input $W_i$ unless that value is explicitly part of the party view.

For finite input and output spaces $\mathcal{X}$ and $\mathcal{O}$ and a size bound $s$, $\C_{\mathcal{X} \to \mathcal{O}, s}$ denotes the class of non-uniform randomized circuits of size at most $s$ that take inputs in $\mathcal{X}$ and output values in $\mathcal{O}$. We use this typed family both for party-view predictors, whose randomized output is an element of $R$, and for Boolean distinguishers used in reduction arguments. When the input/output spaces and the size bound are clear from context, we write $\C$.
  \begin{enumerate}[label=6.\arabic*.]
  \item In the usual complexity-theoretic correspondence, a probabilistic polynomial-time adversary can be represented, for each input length, by a polynomial-size randomized circuit. Conversely, allowing arbitrary polynomial-size circuit families gives the standard non-uniform relaxation of polynomial-time computation \cite[Lecture~4, Thm.~16]{trevisan-randomized-circuits}\cite[Lecture~8]{goldreich-cc-notes}\cite[Chap.~6]{arora-barak-circuits}. Thus, in this note, circuit size is the formal proxy for efficient computation: circuits are finite syntactic objects, and their size can be bounded directly, whereas a full probabilistic-polynomial-time machine model is not formalized here.
  \item A predictor $A \in \C_{\Yspace \to R, s_0}$ observes a view $Y : \Omega \to \Yspace$ and outputs a randomized guess for a target $X \in R$. This is the adversary appearing in the secrecy claim.
  \item A Boolean distinguisher $D \in \C_{\mathcal{Z} \to \{0,1\}, s}$ observes a candidate view or hybrid sample in a finite sample space $\mathcal{Z}$ and outputs a bit in $\{0,1\}$. In the real-or-zero reductions below, such a distinguisher is a proof-side algorithm, not an extra component controlled by Alice: it may receive auxiliary input used to assemble and score a hybrid experiment, even when that auxiliary input is not part of the corrupted party's view.
  \item For any finite view space $\Yspace$, the randomized circuit $A$ induces a Markov kernel $A : \Yspace \to \Delta(R)$, where $\Delta(R)$ is the set of probability distributions on $R$. When presented with a specific view value $y \in \Yspace$, the adversary samples a guess from the distribution $A(y)$.
    \begin{enumerate}[label=6.3.\arabic*.]
    \item In other words, all information observed by the adversary is captured in a view value $y \in \Yspace$. Their strategy $A$ then translates this accumulated knowledge into a biased distribution $A(y) \in \Delta(R)$, from which they sample their guess for the secret value in $R$.
    \end{enumerate}
  \item The average chance that the guess matches the target:
    \[ \Pr[A(Y) = X] := \mathbb{E}_\omega \big[ A(Y(\omega))(\{X(\omega)\}) \big]. \]
  \item The notation $\C_{\mathcal{X} \to \mathcal{O},s}$ is local to this note. The size parameter matters because the hybrid proof turns a predictor $A$ of size at most $s_0$ into a Boolean distinguisher $D_A$ of size at most $s_0 + t_{\mathrm{hyb}}$, where $t_{\mathrm{hyb}}$ accounts for hardwiring the real-or-zero hybrid context, sampling internal randomness, forming the candidate party view from the received pair $(W,C)$ by inserting $C$ into the missing ciphertext position, and comparing $A$'s output with the target value. Thus the real-or-zero assumption must hold for the corresponding typed distinguisher class at the larger size $s_0 + t_{\mathrm{hyb}}$. This follows the resource-bounded viewpoint used in computational entropy \cite[Sec.~2]{hsiao-lu-reyzin} and in computational indistinguishability arguments \cite[Secs.~1.3.3 and 3.2.4]{goldreich-foc}.
    \begin{enumerate}[label=6.5.\arabic*.]
    \item In other words, $D_A$ is a fixed external wrapper around $A$: it contains the circuit for $A$ plus the fixed logic that expands $(W,C)$ into the relevant party view and tests whether $A$ guessed the target. Values appearing only in $W$ may be used by the wrapper for this bookkeeping and scoring, but they are not automatically inputs to the inner predictor $A$; $A$ sees only the party view that $D_A$ constructs. The circuit does not grow during sampling; its larger size comes from this wrapper construction.
    \end{enumerate}
  \end{enumerate}

\item $\HunpC(X \mid Y) := -\log \sup_{A \in \C} \Pr[A(Y) = X]$ is conditional unpredictability entropy \cite[Definition~7]{hsiao-lu-reyzin}. All logarithms are base~$2$.

\item \textbf{Real-or-zero indistinguishability.} $\epscpa \in [0, 1]$ is the distinguishing-probability gap allowed for the fixed public-key registry against the typed distinguisher classes described above. This is the real-or-zero form of the standard hybrid/game-hopping use of IND-CPA security: one ciphertext is replaced at a time, and each replacement is what Shoup calls a transition based on indistinguishability \cite[Sec.~1.1]{shoup-games}.
  \begin{enumerate}[label=8.\arabic*.]
  \item Let $(W,M)$ be efficiently samplable, where $M$ is the challenge plaintext and $W$ is auxiliary input, i.e.\ side information given to the distinguisher. The auxiliary input may be arbitrarily correlated with $M$, and may even contain enough information to recover $M$, as long as it does not contain the matching decryption key $\Dk_p$ or any information that permits decryption under $\pk_p$. Equivalently, even if the auxiliary input $W$ contains the challenge plaintext $M$ itself, or other sampled plaintexts such as $V_2,V_3$, every allowed Boolean distinguisher still has distinguishing advantage at most $\epscpa$, provided $W$ omits the matching decryption key $\Dk_p$ and any decryption-equivalent information for $\pk_p$. For any party $p \in \{a, b, c\}$, the adversary cannot effectively distinguish between the following two views.
    \begin{enumerate}[label=8.1.\arabic*.]
    \item \textbf{Real view:} $(W, \Encf(\pk_p, M; r))$. The adversary sees the auxiliary input $W$ and a fresh encryption of $M$.
    \item \textbf{Zero view:} $(W, \Encf(\pk_p, 0_R; r'))$. The adversary sees the same $W$, but the message is replaced by $0_R$ and encrypted with fresh randomness.
    \end{enumerate}
  \item Formally:
    \[ (W, \Encf(\pk_p, M; r)) \approx_{\C, \epscpa} (W, \Encf(\pk_p, 0_R; r')). \]
  \item \textbf{Key constraints.}
    \begin{enumerate}[label=8.3.\arabic*.]
    \item \textbf{Independence:} the randomness $r, r'$ must be uniform and independent of the message $M$ and side information $W$.
    \item \textbf{Matching-key degeneration:} if $W$ contains the decryption key $\Dk_p$ corresponding to $\pk_p$, then $\Encf(\pk_p, M; r)$ can be decrypted and the real view exposes $M$. Thus the real-or-zero bound is applied only when $W$ omits this matching decryption key and omits any other information that permits decryption under $\pk_p$.
    \item \textbf{Nonmatching keys:} secret keys for other public keys may appear in $W$ only when they do not decrypt ciphertexts under $\pk_p$ and do not reveal $\Dk_p$.
    \end{enumerate}
  \item Equivalently, replacing a fresh encryption of $M$ by a fresh encryption of $0_R$ changes any Boolean distinguishing event by at most $\epscpa$. We use this statement with auxiliary input $W$ in the hybrids below.
  \item \textbf{Soundness as a variant of IND-CPA.} The real-or-zero formulation with auxiliary input is implied by textbook IND-CPA via standard reductions, in three steps.
    \begin{enumerate}[label=8.5.\arabic*.]
    \item \textbf{Real-or-zero $\equiv$ left-or-right $\equiv$ real-or-random.} The three forms of CPA-secure indistinguishability are equivalent up to a factor of $2$ in the distinguishing advantage \cite[Sec.~3]{bellare-concrete}. The real-or-zero form is the convenient hybrid-game-hopping idiom of \cite[Sec.~1.1]{shoup-games}.
    \item \textbf{Auxiliary input.} Indistinguishability with auxiliary input is the standard formulation of semantic security used in uniform-complexity treatments. Goldreich's augmented definition of semantic security explicitly conditions on auxiliary partial information $h(X_n)$ about the message \cite[Def.~4]{goldreich-uniform}, and Goldreich proves it equivalent to indistinguishability of encryptions in the auxiliary-input form \cite[Thm.~1]{goldreich-uniform}. In the non-uniform circuit model used here (Setup item~6), the auxiliary-input formulation is moreover equivalent to the fixed-message formulation, since ``the a-priori information can be incorporated into the circuit'' \cite[Rem.~12]{goldreich-uniform}. The auxiliary input's correlation with $M$ therefore does not impose a stronger assumption.
    \item \textbf{Plaintext-revealing auxiliary input.} Goldreich's indistinguishability definition with auxiliary input $Z_n$ explicitly admits the special case $Z_n = X_n Y_n$, where the auxiliary input contains the challenge plaintexts in cleartext, and ``restricting attention to this special case \ldots\ is equivalent to the general case'' \cite[Def.~5 and Rem.~11]{goldreich-uniform}. The constraint that $W$ omits $\Dk_p$ and any decryption-equivalent information for $\pk_p$ is the standard IND-CPA secrecy condition (the challenger does not reveal the private key to the reduction); without it, no security claim is meaningful. Together, allowing $W$ to contain $M$ in plaintext is the maximally-permissive case of item~8.5.2 and is implied by plain IND-CPA in the non-uniform model.
    \end{enumerate}
  \end{enumerate}

\item \textbf{Notation for indistinguishability.} Following the standard indistinguishability-test formulation \cite[Sec.~3.2]{goldreich-uniform}, $X \approx_{\C, \epscpa} Y$ means that for every Boolean distinguisher $D$ in the relevant typed class,
\[
\left|\Pr[D(X) = 1] - \Pr[D(Y) = 1]\right| \leq \epscpa,
\]
i.e.\ no adversary in $\C$ distinguishes the two distributions with probability gap above $\epscpa$.
\end{enumerate}

\section*{Theorem 1: DSDP secrecy of $V_2$ against corrupted Alice}

\subsection*{Hypotheses}

\begin{enumerate}
\item $V_1, V_2, V_3 : \Omega \to R$ are jointly independent. $V_2$ and $V_3$ are individually uniform on $R$. $V_1$ may have arbitrary distribution; only its independence from $(V_2, V_3)$ is used in the proof below.

\item $U_1, U_2, U_3, R_2, R_3 : \Omega \to R$ satisfy the following properties, which are the ones used in the proof of this theorem:
  \begin{enumerate}[label=2.\arabic*.]
  \item $(V_1, V_2, V_3) \indep (U_1, U_2, U_3)$ and the $U_i$ are mutually independent.
  \item $U_3 \in R^\times$ almost surely. That is, $\Pr[U_3 \notin R^\times] = 0$, so $U_3$ is neither $0_R$ nor a zero divisor with positive probability, a case that belongs to the malicious-adversary analysis. Its distribution on $R^\times$ is otherwise arbitrary.
  \item $(R_2, R_3) \indep (V_1, V_2, V_3, U_1, U_2, U_3)$.
  \end{enumerate}

  \textit{Remark.} At the protocol level, $U_1, U_2$ and $R_2, R_3$ are uniform on $R$, since these uniformities are required by the analogous secrecy theorems against corrupted Bob and corrupted Charlie. The current theorem does not invoke them, so they are omitted from the hypotheses above. A reader composing this theorem with the corrupted-Bob and corrupted-Charlie analogues must reinstate them in the global hypothesis set.

\item Protocol-defined intermediate random variables on $\Omega$:
\begin{align*}
D_2 &:= V_2 \cdot U_2 + R_2, \\
D_3 &:= V_3 \cdot U_3 + R_3 + D_2 = V_3 \cdot U_3 + V_2 \cdot U_2 + R_2 + R_3, \\
S   &:= D_3 + U_1 \cdot V_1 - R_2 - R_3 = U_1 \cdot V_1 + U_2 \cdot V_2 + U_3 \cdot V_3,
\end{align*}
where $S$ is the publicly opened sum at the reveal step.

\item $r_a, r_b, r_c : \Omega \to \Renc$ are mutually independent, each individually uniform on $\Renc$, and jointly independent of $(V_1, V_2, V_3, U_1, U_2, U_3, R_2, R_3)$. Namely,
\[
(V_1, V_2, V_3, U_1, U_2, U_3, R_2, R_3) \indep r_a \indep r_b \indep r_c.
\]

\item The corrupted-Alice view is
\begin{multline*}
\AliceView := \big( \Dk_a, S, V_1, U_1, U_2, U_3, R_2, R_3, \\
\Encf(\pk_a, D_3; r_a), \Encf(\pk_c, V_3; r_c), \Encf(\pk_b, V_2; r_b) \big).
\end{multline*}
\end{enumerate}

\subsection*{Conclusion}

For every allowed predictor $A$ of size at most $s_0$,
\[ \Pr[A(\AliceView) = V_2] \leq \tfrac{1}{m} + 2 \cdot \epscpa, \]
and therefore
\[
\HunpC(V_2 \mid \AliceView)
\geq -\log\left(\frac{1}{m} + 2 \epscpa\right)
= \log m - \log(1 + 2m\epscpa).
\]

\subsection*{Proof outline}

The bound is established by a hybrid argument. We apply real-or-zero indistinguishability twice, replacing each cross-key ciphertext by an encryption of $0_R$, and then use an information-theoretic uniformity argument on the resulting hybrid.

\textit{Probability-space extension.} The proof uses two fresh randomness variables $r'_c, r'_b : \Omega \to \Renc$ that are not part of the protocol. We take $r'_c$ uniform on $\Renc$ and jointly independent of $(V_1,V_2,V_3,U_1,U_2,U_3,R_2,R_3,r_a,r_b,r_c)$, and $r'_b$ uniform on $\Renc$ and jointly independent of $(V_1,V_2,V_3,U_1,U_2,U_3,R_2,R_3,r_a,r_b,r_c,r'_c)$.

\textit{Views used in the proof.} Recall from Hypothesis~5 that the corrupted-Alice view is
\begin{multline*}
\AliceView = \big( \Dk_a, S, V_1, U_1, U_2, U_3, R_2, R_3, \\
\Encf(\pk_a, D_3; r_a), \Encf(\pk_c, V_3; r_c), \Encf(\pk_b, V_2; r_b) \big).
\end{multline*}
The two intermediate hybrid views are
\begin{multline*}
H_1 := \big( \Dk_a, S, V_1, U_1, U_2, U_3, R_2, R_3, \\
\Encf(\pk_a, D_3; r_a), \Encf(\pk_c, 0_R; r'_c), \Encf(\pk_b, V_2; r_b) \big),
\end{multline*}
\begin{multline*}
H_2 := \big( \Dk_a, S, V_1, U_1, U_2, U_3, R_2, R_3, \\
\Encf(\pk_a, D_3; r_a), \Encf(\pk_c, 0_R; r'_c), \Encf(\pk_b, 0_R; r'_b) \big).
\end{multline*}
The view used for the residual uniformity argument extends $H_2$ by the decrypted plaintext $D_3$:
\begin{multline*}
\AliceViewHyb := \big( \Dk_a, S, V_1, U_1, U_2, U_3, R_2, R_3, D_3, \\
\Encf(\pk_a, D_3; r_a), \Encf(\pk_c, 0_R; r'_c), \Encf(\pk_b, 0_R; r'_b) \big).
\end{multline*}
The two reduction-side auxiliary inputs are
\begin{multline*}
W_1 := \big( V_2, V_3, \Dk_a, S, V_1, U_1, U_2, U_3, R_2, R_3, \\
\Encf(\pk_a, D_3; r_a), \Encf(\pk_b, V_2; r_b) \big),
\end{multline*}
\begin{multline*}
W_2 := \big( V_2, V_3, \Dk_a, S, V_1, U_1, U_2, U_3, R_2, R_3, \\
\Encf(\pk_a, D_3; r_a), \Encf(\pk_c, 0_R; r'_c) \big).
\end{multline*}
\textit{Proof chain.} Steps~1 to~3 relate $\AliceView$ to $\AliceViewHyb$ through two computational indistinguishability transitions and one deterministic extension:
\[
\AliceView \;\xRightarrow[\text{external }D_A\text{ with }W_1]{\text{Step 1}}\; H_1
\;\xRightarrow[\text{external }D_A\text{ with }W_2]{\text{Step 2}}\; H_2
\;\xrightarrow{\text{Step 3: }D_3\text{-decryption}}\; \AliceViewHyb.
\]
The double arrows record proof-side indistinguishability arguments, not computations available to corrupted Alice. Setup item~6 separates the party-view predictor $A$ from the external wrapper distinguisher $D_A$, and Setup item~8 is used in its auxiliary-input form. Thus, in Steps~1 and~2, $D_A$ may receive the corresponding $W_i$ and use the sampled plaintexts stored there for bookkeeping and scoring, while the predictor $A$ inside the wrapper is given only the candidate Alice view. Each of Steps~1 and~2 replaces one cross-key ciphertext by an encryption of $0_R$ and pays $\epscpa$ in success-probability gap; the auxiliary inputs $W_1$ and $W_2$ are the data that the corresponding reduction shares between the real and zero sides of the swap. Step~3 decrypts Alice's own-key ciphertext to recover the plaintext $D_3$ and adds it to $H_2$ as an explicit component (zero cost). Step~4 then bounds $\Pr[V_2 = v \mid \AliceViewHyb] = 1/m$.

Relating each view to $\AliceView$:
\begin{itemize}
\item $H_1$ replaces $\Encf(\pk_c, V_3; r_c)$ in $\AliceView$ with $\Encf(\pk_c, 0_R; r'_c)$.
\item $H_2$ further replaces $\Encf(\pk_b, V_2; r_b)$ in $H_1$ with $\Encf(\pk_b, 0_R; r'_b)$.
\item $\AliceViewHyb$ extends $H_2$ by the decrypted plaintext $D_3$, recovered via $\Dk_a$ from the own-key ciphertext.
\item $W_1$ is $\AliceView$ with $\Encf(\pk_c, V_3; r_c)$ removed and $V_2, V_3$ added as auxiliary plaintexts for the Step~1 reduction.
\item $W_2$ is $H_1$ with $\Encf(\pk_b, V_2; r_b)$ removed and $V_2, V_3$ retained as auxiliary plaintexts for the Step~2 reduction.
\end{itemize}

\paragraph{Step 1, first ciphertext swap.} The challenge key is $\pk_c$ and the challenge message is $V_3$. The auxiliary input $W_1$ contains the sampled plaintexts $V_2,V_3$, but it does not contain the matching decryption key $\Dk_c$ or any information that permits decryption under $\pk_c$. Therefore the real-or-zero assumption applies with $W_1$, target key $\pk_c$, and message $V_3$:
\[
(W_1, \Encf(\pk_c, V_3; r_c)) \approx_{\C, \epscpa} (W_1, \Encf(\pk_c, 0_R; r'_c)).
\]
By the probability-space extension, $r'_c$ is jointly independent of $(V_1,V_2,V_3,U_1,U_2,U_3,R_2,R_3,r_a,r_b,r_c)$, and since $W_1$ and $V_3$ are deterministic functions of these variables, $r'_c$ is in particular jointly independent of $(W_1, V_3, r_c)$.

For any allowed predictor $A$ of size at most $s_0$, form the Boolean distinguisher $D_A$ that receives a pair $(W_1,C)$, forms the corresponding Alice view by inserting $C$ into the missing $\pk_c$ ciphertext position, samples $A$'s guess on that view, and outputs $1$ iff the guess equals the auxiliary target value $V_2$. The auxiliary-input assumption is being used in its stronger form here: even if the external wrapper is given the sampled plaintexts in $W_1$, the predictor $A$ inside the wrapper still receives only the candidate party view, and the real-or-zero gap is bounded by $\epscpa$. By the wrapper-overhead convention in Setup item~6, $D_A$ belongs to the typed distinguisher class covered by the real-or-zero assumption. By construction,
\[
\Pr[D_A(W_1,\Encf(\pk_c,V_3;r_c))=1] = \Pr[A(\AliceView)=V_2]
\]
and
\[
\Pr[D_A(W_1,\Encf(\pk_c,0_R;r'_c))=1] = \Pr[A(H_1)=V_2].
\]
Thus the distinguishing bound for $D_A$ gives
\[
\left|\Pr[A(\AliceView)=V_2] - \Pr[A(H_1)=V_2]\right| \leq \epscpa.
\]

\paragraph{Step 2, second ciphertext swap.} The challenge key is $\pk_b$ and the challenge message is $V_2$. The auxiliary input $W_2$ contains the sampled plaintexts $V_2,V_3$, but it does not contain the matching decryption key $\Dk_b$ or any information that permits decryption under $\pk_b$. Therefore real-or-zero indistinguishability applies with $W_2$, target key $\pk_b$, and message $V_2$. This is a known-message use of real-or-zero: the auxiliary input $W_2$ explicitly contains $V_2$, which is allowed by Setup item~8 because the matching decryption key is absent.
\[
(W_2, \Encf(\pk_b, V_2; r_b)) \approx_{\C, \epscpa} (W_2, \Encf(\pk_b, 0_R; r'_b)).
\]
By the probability-space extension, $r'_b$ is jointly independent of $(V_1,V_2,V_3,U_1,U_2,U_3,R_2,R_3,r_a,r_b,r_c,r'_c)$, and since $W_2$ and $V_2$ are deterministic functions of these variables, $r'_b$ is in particular jointly independent of $(W_2, V_2, r_b)$.

The same predictor-to-distinguisher conversion forms a wrapper $D_A$ that receives $(W_2,C)$, inserts $C$ into the missing $\pk_b$ ciphertext position, runs $A$ only on the resulting candidate view, and compares the guess with the auxiliary target $V_2$ inside the wrapper. The auxiliary-input assumption is again being used in its stronger form: even if the external wrapper is given the sampled plaintexts retained in $W_2$, the inner predictor $A$ does not receive them as plaintext components, and the real-or-zero gap is bounded by $\epscpa$. Using the wrapper-overhead convention from Setup item~6 gives
\[
\left|\Pr[A(H_1)=V_2] - \Pr[A(H_2)=V_2]\right| \leq \epscpa.
\]

\paragraph{Step 3, own-key decryption.} The ciphertext $\Encf(\pk_a, D_3; r_a)$ is encrypted under Alice's own public key, and $\Dk_a$ is in corrupted Alice's view. By correctness, $D_3$ is a deterministic function of $(\Dk_a,\Encf(\pk_a,D_3;r_a))$. Recall the view $\AliceViewHyb$ from the proof outline, which extends $H_2$ by the explicit component $D_3$. Any size-$s_0$ predictor $A_2$ on $H_2$ extends to a size-$(s_0 + O(1))$ predictor $A'_2$ on $\AliceViewHyb$ that ignores the explicit $D_3$ component and feeds the remaining components to $A_2$. The two predictors achieve identical success probability:
\[
\Pr[A_2(H_2) = V_2] = \Pr[A'_2(\AliceViewHyb) = V_2].
\]
Hence it suffices to bound $\Pr[A'(\AliceViewHyb) = V_2]$ for every size-$(s_0 + O(1))$ predictor $A'$ on $\AliceViewHyb$. Given $D_3$ and $\Dk_a$, the own-key ciphertext is a deterministic function of the fresh randomness $r_a$, which is independent of $V_2$ and of the other input variables. Hence it carries no additional information about $V_2$ once $D_3$ is conditioned on.

\paragraph{Step 4, residual uniformity.}

\textbf{Initial goal.} Show that $\Pr[V_2 = v \mid \AliceViewHyb] = 1/m$ for every $v \in R$ and every positive-probability conditioning value of $\AliceViewHyb$, where $\AliceViewHyb$ is defined in the proof outline. We reach the goal by reducing the conditioning in four steps.

\medskip

\textit{Substep 4.1, drop the residual encryptions.} Set
\begin{align*}
Y &:= (\Dk_a, S, V_1, U_1, U_2, U_3, R_2, R_3, D_3), \\
W &:= (r_a, r'_c, r'_b), \\
Z &:= \big(\Encf(\pk_a, D_3; r_a),\ \Encf(\pk_c, 0_R; r'_c),\ \Encf(\pk_b, 0_R; r'_b)\big),
\end{align*}
so that $\AliceViewHyb = (Y, Z)$ and $Z$ is a deterministic function of $(Y, W)$, obtained by reading $D_3$ out of $Y$ and combining it with the public constants $\pk_a, \pk_c, \pk_b, 0_R$ and the components of $W$. By Hypothesis~4 (for $r_a$) and the probability-space extension in the proof outline (for $r'_c$ and $r'_b$), $W \indep (V_2, Y)$. By \texttt{inde\_RV2\_cinde} from infotheo (in \texttt{du2002/spp\_proba.v}), this yields the conditional independence $V_2 \indep W \mid Y$, and hence $V_2 \indep Z \mid Y$ since $Z$ is a function of $(Y, W)$. By \texttt{cinde\_rv\_comp\_removal} from infotheo (in \texttt{du2002/spp\_entropy.v}), this implies
\[
\Pr[V_2 = v \mid \AliceViewHyb] = \Pr[V_2 = v \mid Y]
\]
for every $v \in R$ and every positive-probability conditioning value of $\AliceViewHyb$.

\quad Goal becomes: $\Pr[V_2 = v \mid (\Dk_a, S, V_1, U_1, U_2, U_3, R_2, R_3, D_3)] = 1/m$ for every $v \in R$.

\medskip

\textit{Substep 4.2, drop $D_3$ and $\Dk_a$.} $D_3 = (S - U_1 V_1) + R_2 + R_3$ is a deterministic function of $(S, V_1, U_1, R_2, R_3)$, all already in the conditioning. Conditioning on a deterministic function of variables already in the conditioning does not change the conditional distribution. $\Dk_a$ is a deterministic constant fixed at protocol setup, independent of all input random variables, so it also does not affect the conditional distribution.

\quad Goal becomes: $\Pr[V_2 = v \mid (S, V_1, U_1, U_2, U_3, R_2, R_3)] = 1/m$ for every $v \in R$.

\medskip

\textit{Substep 4.3, drop $R_2, R_3$.} By hypothesis~2,
\[
(R_2, R_3) \indep (V_1, V_2, V_3, U_1, U_2, U_3).
\]
The opened sum $S = U_1 V_1 + U_2 V_2 + U_3 V_3$ is a function of $(V_1, V_2, V_3, U_1, U_2, U_3)$, so $(R_2, R_3)$ is jointly independent of the whole tuple $(S, V_1, U_1, U_2, U_3)$. Conditioning on these independent variables does not change the conditional distribution.

\quad Goal becomes: $\Pr[V_2 = v \mid (S, V_1, U_1, U_2, U_3)] = 1/m$ for every $v \in R$.

\medskip

\textit{Substep 4.4, solve the linear equation.} Set $c := S - U_1 V_1$, a constant given the conditioning on $(S, V_1, U_1)$. The remaining equation is $U_2 V_2 + U_3 V_3 = c$ with unknowns $V_2$ and $V_3$. For each candidate $v \in R$ for $V_2$, the unique $V_3$ satisfying the equation is $V_3 = U_3^{-1}(c - U_2 v)$, well-defined because $U_3 \in R^\times$ by hypothesis~2. So the fiber of solutions has exactly $m$ elements, one per candidate $v$. By Hypothesis~1, $(V_2, V_3)$ is jointly uniform on $R \times R$. By Hypothesis~2.1, $(V_2, V_3) \indep (V_1, U_1, U_2, U_3)$, so conditioning on $(V_1, U_1, U_2, U_3)$ does not change the joint law of $(V_2, V_3)$, which remains uniform on $R \times R$. The constraint $U_2 V_2 + U_3 V_3 = c$ then partitions $R \times R$ into $m$ fibers of size $m$, and the conditional law of $(V_2, V_3)$ given the equation is uniform on the relevant fiber. Marginalizing $V_3$ gives $V_2$ uniform on $R$.

\quad Goal achieved: $\Pr[V_2 = v \mid (S, V_1, U_1, U_2, U_3)] = 1/m$ for every $v \in R$.

This is the abstract finite-ring version of the fiber argument formalized as \texttt{Pr\_dsdp\_sol\_uniform} in the existing infotheo development. Combined with substeps~4.1, 4.2, 4.3, the initial goal $\Pr[V_2 = v \mid \AliceViewHyb] = 1/m$ for every $v \in R$ and every positive-probability conditioning value is established.

\paragraph{Step 5, conclude.} Since $D_3$ is computable from $H_2$, any predictor on $H_2$ can be viewed as a predictor on $\AliceViewHyb$ that ignores the explicit $D_3$ component. By Step~4, this predictor succeeds with probability exactly $1/m$:
\[
\Pr[A(H_2) = V_2] = \frac{1}{m} \quad \text{for every allowed predictor } A.
\]
Steps~1 and~2 establish
\begin{align*}
\big|\Pr[A(\AliceView)=V_2] - \Pr[A(H_1)=V_2]\big| &\leq \epscpa, \\
\big|\Pr[A(H_1)=V_2] - \Pr[A(H_2)=V_2]\big| &\leq \epscpa.
\end{align*}
In particular, $\Pr[A(\AliceView)=V_2] - \Pr[A(H_1)=V_2] \leq \epscpa$ and $\Pr[A(H_1)=V_2] - \Pr[A(H_2)=V_2] \leq \epscpa$. Adding the two inequalities gives
\begin{align*}
\Pr[A(\AliceView) = V_2] &\leq \Pr[A(H_2) = V_2] + 2 \cdot \epscpa \\
                          &= \tfrac{1}{m} + 2 \cdot \epscpa.
\end{align*}
The entropy form is derived from the probability bound in three steps.

\textit{Step 5.1, take the supremum over allowed predictors.} The bound $\Pr[A(\AliceView) = V_2] \leq \frac{1}{m} + 2 \epscpa$ holds for every allowed predictor $A$ of size at most $s_0$. Since the right-hand side $\frac{1}{m} + 2 \epscpa$ is the same number for every choice of $A$, taking the supremum over $A$ on the left preserves the inequality:
\[
\sup_{A \in \C} \Pr[A(\AliceView) = V_2] \leq \frac{1}{m} + 2 \epscpa.
\]

\textit{Step 5.2, apply $-\log$.} The function $x \mapsto -\log x$ is monotonically decreasing on $(0, \infty)$. Applying it to both sides reverses the inequality:
\[
-\log \sup_{A \in \C} \Pr[A(\AliceView) = V_2] \;\geq\; -\log\left(\frac{1}{m} + 2 \epscpa\right).
\]
By the definition of conditional unpredictability entropy in Setup item~7, the left-hand side is exactly $\HunpC(V_2 \mid \AliceView)$, hence
\[
\HunpC(V_2 \mid \AliceView) \;\geq\; -\log\left(\frac{1}{m} + 2 \epscpa\right).
\]

\textit{Step 5.3, simplify the right-hand side.} Factor $1/m$ inside the logarithm and split the product:
\begin{align*}
-\log\left(\frac{1}{m} + 2 \epscpa\right)
&= -\log\left(\frac{1}{m} \cdot \big(1 + 2\,m\,\epscpa\big)\right) \\
&= -\log\frac{1}{m} \,-\, \log\!\big(1 + 2\,m\,\epscpa\big) \\
&= \log m \,-\, \log\!\big(1 + 2\,m\,\epscpa\big).
\end{align*}
Combining with Step~5.2 gives the stated entropy bound:
\[
\HunpC(V_2 \mid \AliceView) \;\geq\; \log m \,-\, \log\!\big(1 + 2\,m\,\epscpa\big).
\]
This bound is informative when $\frac{1}{m} + 2\,\epscpa < 1$. Otherwise it remains formally true but weak.

\section*{Reference Notes}

Reference \cite[Definition~7]{hsiao-lu-reyzin} introduces conditional unpredictability entropy with the adversary class fixed concretely as polynomial-size circuits. It also states in Section~2: ``We will use $s$ as the circuit size parameter (or running time bound when dealing with Turing machines instead of circuits).'' The typed notation for $\C_{\mathcal{X} \to \mathcal{O},s}$ is local to this 3-party Alice note; its role is to record the resource bound needed when a predictor is wrapped into a larger hybrid distinguisher. Reference \cite[Secs.~1.3.3 and 3.2.4]{goldreich-foc} is a textbook source for the non-uniform circuit model and indistinguishability by circuits.

Reference \cite[Sec.~2.1]{lindell-pinkas-yao} is a standard source for the static semi-honest setting: such an adversary controls one party, follows the protocol specification, and tries to learn extra information from the transcript/messages in that party's view. The external Boolean distinguisher used in the hybrids above is reduction machinery around this party-view adversary, not additional protocol information given to the corrupted party.

For the bridge from probabilistic polynomial-time computation to polynomial-size circuits, Trevisan's lecture notes introduce randomized algorithms and state that they ``show that Boolean circuits can simulate randomized algorithms'' \cite[Lecture~4]{trevisan-randomized-circuits}; the same notes state Adleman's theorem as $\mathrm{BPP} \subseteq \mathrm{SIZE}(n^{O(1)})$ and conclude that one can simulate the probabilistic algorithm by ``a circuit of size polynomial in $n$'' \cite[Lecture~4, Thm.~16]{trevisan-randomized-circuits}. Goldreich's computational complexity notes describe $P/\mathrm{poly}$ as a non-uniform model and say that $P/\mathrm{poly}$ ``upper bounds'' efficient computation, with $\mathrm{BPP} \subseteq P/\mathrm{poly}$ \cite[Lecture~8]{goldreich-cc-notes}. Arora and Barak define $P/\mathrm{poly}$ as languages decidable by polynomial-sized circuit families and state $P \subseteq P/\mathrm{poly}$, while also recording that $P$-uniform circuit families characterize $P$ \cite[Chap.~6]{arora-barak-circuits}.

References \cite{bellare-concrete,shoup-games,blanchet-pointcheval-games} justify the proof style used above: encryption indistinguishability is commonly handled through real-or-random and left-or-right style notions, and security proofs are often organized as game-hopping or hybrid sequences whose adjacent transitions are justified by computational indistinguishability. The public-key real-or-zero security of the fixed registry is a hypothesis of this note. The proof style is instantiated by two cross-key ciphertext swaps followed by the information-theoretic residual uniformity argument corresponding to \texttt{Pr\_dsdp\_sol\_uniform}.

% =============================================================================
% Verbatim quotes from cited references supporting Setup items 8.5.2 and 8.5.3.
% Located here for traceability; not rendered in the typeset document.
% =============================================================================
%
% --- Supporting Setup item 8.5.2 (auxiliary input is the standard formulation
% --- in uniform-complexity treatments, and equivalent to fixed-message
% --- formulation in the non-uniform circuit model). Source:
% --- Goldreich, "A Uniform-Complexity Treatment of Encryption and
% --- Zero-Knowledge", Section 3 ("Uniform Security of Encryption Schemes").
%
% Quote 1 (Section 3.1, motivation paragraph before Definition 4):
% "Loosely speaking, semantic security as defined in [18] means that whatever
%  can be efficiently computed from the encryption of the message, can be
%  efficiently computed given only the length of the message. Here we augment
%  this definition by requiring that the above remains valid in presence of
%  auxiliary partial information about the message. Namely, whatever can be
%  efficiently computed from the encryption of the message and additional
%  partial information about the message, can be efficiently computed given
%  only the length of the message and the same partial information."
%
% Quote 2 (Definition 4):
% "An encryption scheme, (G, E, D), is semantically secure if for every
%  probabilistic polynomial-time algorithm A, there exists a probabilistic
%  polynomial-time algorithm A', such that for every polynomial-time random
%  variable {X_n}, every polynomial-time computable function h, every function
%  f, every constant c > 0 and all sufficiently large n:
%  Pr[A(E(X_n), h(X_n), 1^n) = f(X_n)]
%    < Pr[A'(|X_n|, h(X_n), 1^n) = f(X_n)] + 1/n^c."
% (h(X_n) is the auxiliary input correlated with the message.)
%
% Quote 3 (Theorem 1):
% "An encryption scheme is semantically secure if and only if it has
%  indistinguishable encryptions."
%
% Quote 4 (Remark 12, non-uniform equivalence -- the key load-bearing quote):
% "In the non-uniform case, the definition is equivalent to requiring that,
%  for all families of polynomial-size circuits {C_n}, all sequences of pairs
%  {(alpha_n, beta_n)} (with |alpha_n| = |beta_n|), and all c > 0 and all
%  sufficiently large n:
%  | Pr[C_n(E(alpha_n)) = 1] - Pr[C_n(E(beta_n)) = 1] | < 1/n^c.
%  (Namely, in this case the a-priori information can be incorporated into
%  the circuit.)"
%
% --- Supporting Setup item 8.5.3 (auxiliary input may contain the challenge
% --- plaintexts in cleartext; the maximally permissive case is equivalent to
% --- the general case). Source: same Goldreich paper, Section 3.2.
%
% Quote 5 (Definition 5, indistinguishability with auxiliary input):
% "An encryption scheme, (G, E, D), has indistinguishable encryptions if for
%  every polynomial-time random variable {T_n = X_n Y_n Z_n} (with
%  |X_n| = |Y_n|), every probabilistic polynomial-time algorithm A, every
%  constant c > 0 and all sufficiently large n:
%  | Pr[A(Z_n, E(X_n)) = 1] - Pr[A(Z_n, E(Y_n)) = 1] | < 1/n^c."
% (Z_n is the auxiliary input jointly distributed with X_n, Y_n.)
%
% Quote 6 (Remark 11 -- direct support for the plaintext-revealing case):
% "The random variable Z_n models additional information, on the message
%  space, given to algorithm A which tries to distinguish the encryptions
%  (of X_n and Y_n). A special case of interest is when Z_n = X_n Y_n. In view
%  of our results (see subsection 3.3), restricting attention to this special
%  case (as done in [18]) is equivalent to the general case."
%
% Quote 7 (Proof of Proposition 1 -- explicit use of plaintext-containing
% auxiliary input in the construction):
% "we can distinguish the encryptions of X_n and Y_n =def 1^|X_n| (using
%  auxiliary input Z_n = X_n Y_n)."
%
% --- Note: the Goldreich paper does NOT explicitly state the "W omits Dk_p"
% --- constraint. That constraint is the standard IND-CPA secrecy condition
% --- (the challenger does not reveal the private key to the reduction); it
% --- is implicit in any IND-CPA-style security claim and noted in 8.5.3 only
% --- to make the soundness condition for our reduction self-contained.
%
% --- Supporting Setup item 8.5.1 (real-or-zero, real-or-random, and
% --- left-or-right CPA indistinguishability are equivalent up to a small
% --- constant factor; the real-or-zero form is the standard hybrid-game-
% --- hopping idiom). Sources: Bellare-Desai-Jokipii-Rogaway 2000 and
% --- Shoup 2006.
%
% Quote 8 (Bellare-Desai-Jokipii-Rogaway, Section 3, equivalence summary):
% "We show that LOR-ATK and ROR-ATK are equivalent, up to a small constant
%  factor in the reduction. (That is, we have security-preserving reductions
%  between them.) We also show a security-preserving reduction from these
%  notions to FTG-ATK. ... We complete the picture by showing that SEM-ATK
%  and FTG-ATK are equivalent."
%
% Quote 9 (Bellare-Desai-Jokipii-Rogaway, Theorem 1, the factor-of-2 bound):
% "Theorem 1 [ROR-ATK => LOR-ATK] For any scheme SE = (K, E, D),
%   Adv^{lor-cpa}_{SE}(k, t, q_e, mu_e) <= 2 * Adv^{ror-cpa}_{SE}(k, t, q_e, mu_e)."
% (Theorem 2 gives the reverse direction with factor 1, so ROR ~ LOR with
%  worst-case factor 2.)
%
% --- Note: BDJR does not use the term "real-or-zero" explicitly. Real-or-zero
% --- is the special case of LOR where one of the two messages is fixed to
% --- 0^|x|. By BDJR Theorem 1, indistinguishability under any such fixed
% --- LOR oracle inherits the LOR <-> ROR equivalence up to the same constant.
% --- Shoup uses real-or-zero / single-message-substitution hybrids
% --- throughout the games tutorial.
%
% Quote 10 (Shoup, Section 1.1, definition of indistinguishability transition):
% "Transitions based on indistinguishability. In such a transition, a small
%  change is made that, if detected by the adversary, would imply an efficient
%  method of distinguishing between two distributions that are
%  indistinguishable (either statistically or computationally)."
%
% Quote 11 (Shoup, Section 1.1, hybrid-with-auxiliary-input construction --
%  this is the construction we use in Steps 1 and 2):
% "Typically, one designs the two games so that they could easily be rewritten
%  as a single 'hybrid' game that takes an auxiliary input -- if the auxiliary
%  input is drawn from P1, you get Game i, and if drawn from P2, you get
%  Game i + 1. The distinguisher then simply runs this single hybrid game
%  with its input, and outputs 1 if the appropriate event occurs."
%
% Quote 12 (Shoup, Section 1.1, hybrid-argument continuity with the literature):
% "'Hybrid arguments' have been used extensively in cryptography for many
%  years. Such an argument is essentially a sequence of transitions based on
%  indistinguishability."
% =============================================================================

\begin{thebibliography}{13}\bibitem{hsiao-lu-reyzin}
Hsiao, Lu, Reyzin.
\emph{Conditional Computational Entropy, or Toward Separating Pseudoentropy from Compressibility}.
Advances in Cryptology -- Eurocrypt 2007, Lecture Notes in Computer Science 4515, pages 169--186.
\url{https://www.iacr.org/archive/eurocrypt2007/45150169/45150169.pdf}

\bibitem{lindell-pinkas-yao}
Lindell, Pinkas.
\emph{A Proof of Security of Yao's Protocol for Two-Party Computation}.
Journal of Cryptology, Vol.~22, No.~2, pages 161--188, 2009; full version, June 26, 2006.
\url{https://eprint.iacr.org/2004/175.pdf}

\bibitem{bellare-concrete}
Bellare, Desai, Jokipii, Rogaway.
\emph{A Concrete Security Treatment of Symmetric Encryption}.
Full version of FOCS 1997 paper, September 2000.
\url{https://cseweb.ucsd.edu/~mihir/papers/sym-enc.pdf}

\bibitem{shoup-games}
Shoup.
\emph{Sequences of Games: A Tool for Taming Complexity in Security Proofs}.
January 18, 2006.
\url{https://www.shoup.net/papers/games.pdf}

\bibitem{blanchet-pointcheval-games}
Blanchet, Pointcheval.
\emph{Automated Security Proofs with Sequences of Games}.
Advances in Cryptology -- Crypto 2006, Lecture Notes in Computer Science 4117, pages 537--554.
\url{https://bblanche.gitlabpages.inria.fr/publications/BlanchetPointchevalCrypto06.pdf}

\bibitem{goldreich-foc}
Goldreich.
\emph{Foundations of Cryptography, Volume I: Basic Tools}.
Cambridge University Press, 2001.
\url{https://www.wisdom.weizmann.ac.il/~oded/foc-vol1.html}

\bibitem{goldreich-uniform}
Goldreich.
\emph{A Uniform-Complexity Treatment of Encryption and Zero-Knowledge}.
Journal of Cryptology, Vol.~6, No.~1, pages 21--53, 1993; revised version, July 1998.
\url{https://www.wisdom.weizmann.ac.il/~oded/PSX/uniform.pdf}

\bibitem{trevisan-randomized-circuits}
Trevisan.
\emph{CS254: Computational Complexity, Notes for Lecture 4}.
January 16, 2014.
\url{https://theory.stanford.edu/~trevisan/cs254-14/lecture04.pdf}

\bibitem{goldreich-cc-notes}
Goldreich.
\emph{Computational Complexity: A Conceptual Perspective, lecture notes}.
Lecture~8: Non-Uniform Polynomial Time ($P/\mathrm{poly}$).
\url{https://www.wisdom.weizmann.ac.il/~oded/PS/CC/all.pdf}

\bibitem{arora-barak-circuits}
Arora, Barak.
\emph{Computational Complexity: A Modern Approach}, Chapter~6 draft: Boolean Circuits.
\url{https://theory.cs.princeton.edu/complexity/circuitschap.pdf}

\end{thebibliography}

\end{document}
